\section{Related Work}

\label{sec:related-work}

% PAPER_STUDIO_PARAGRAPH:RW-P1

Profile-conditioned dialogue systems condition generation on a static user profile, yet the attributes that govern an e-commerce exchange, such as a shopper's budget, delivery preference, or willingness to accept a substitute, shift over the course of a conversation. A fixed profile snapshot cannot capture these dynamics because it reflects a prior state rather than the evolving interaction. Consequently, such models may persist with outdated assumptions, recommending a high-end item after the user has signaled a tighter budget or failing to revisit a preference that the dialogue itself has contradicted. Personalization on a static representation thus addresses a narrow slice of the problem: it improves how faithfully utterances reflect a stored persona, but it does not provide a mechanism for revising the reasoning that connects new conversational evidence to a changing attribute set. Reasoning over fresh dialogue context is therefore the complementary capability that profile conditioning lacks.

% PAPER_STUDIO_PARAGRAPH:RW-P2

Task-oriented dialogue systems are typically evaluated on their ability to complete a well-specified transaction, such as booking a flight or reserving a restaurant, where success is defined by slot filling and API calls. E-commerce dialogue introduces a distinct set of constraints that this formulation does not capture: the system must reason over product attributes that vary across categories, balance competing objectives such as attribute fidelity and conversational naturalness, and adapt to users who may not articulate their constraints in a single turn. Moreover, the cost of an incorrect recommendation is asymmetric; a fluent but wrong answer can erode user trust more than a correct but awkward one. These properties place e-commerce dialogue at the intersection of task completion and open-ended generation, where neither a strict slot-filling evaluation nor a purely fluency-based metric is sufficient to characterize system quality. Our evaluation therefore combines task-oriented measures with human-aligned response-quality dimensions to reflect this dual requirement.

% PAPER_STUDIO_PARAGRAPH:RW-P3

Reinforcement learning for dialogue systems commonly optimizes a single scalar reward, whether task success, informativeness, or user satisfaction, with the implicit assumption that one aggregated signal sufficiently captures response quality. In e-commerce settings, however, this assumption breaks down because the correctness of the requested attribute and the naturalness of the phrasing can conflict within a single turn, and the relative importance of these dimensions shifts with context. A fixed linear combination of rewards encodes a static trade-off that cannot adapt when a user prioritizes a precise product match over conversational polish, or vice versa. Our approach instead coordinates multiple reward channels by assigning rank-based gradient weights per training batch, so the optimization responds to the current distribution of candidate quality rather than committing to a predetermined weighting. This contrasts with single-reward formulations that cannot express such context-sensitive coordination without manual reward tuning.
